\chapter{Pascal}\label{cha:pascal}

The machine has two Pascals, and they share nothing but the language.
Keep them apart in your head from the start:

\begin{center}
\small
\begin{tabular}{@{}l >{\raggedright\arraybackslash}p{0.33\linewidth} >{\raggedright\arraybackslash}p{0.33\linewidth}@{}}
\toprule
 & \textbf{Turbo Pascal 3} & \textbf{Mad Pascal} \\
\midrule
What it is & a program of 1985 & a cross-compiler of today \\
Where it runs & on the machine: the Z80, under CP/M & on your desktop \\
What it makes & a CP/M \texttt{.COM}, for the Z80 & a K/OS \texttt{.prg}, for the 45GS10 \\
What it can touch & CP/M's world only & every chip in the machine \\
How you get it & you supply it; it is Borland's & in the repository, free \\
\bottomrule
\end{tabular}
\end{center}

\noindent Nothing carries over between them --- not the source, not the
toolchain, not the programs they produce. The first section is one, the
second is the other.

\section{Turbo Pascal 3, on CP/M}
Borland's Turbo Pascal~3 is a CP/M program, so it runs where CP/M runs:
on the Tube's Z80\shipif{chapter:09-cpm}{ (Chapter~\ref{cha:cpm})}, in its own world, exactly as
it did on a Kaypro. You write, compile and run without leaving it, and
what it produces is a CP/M program that runs under CP/M on the Z80 ---
it cannot see VICKY, the OPL2 or anything else on the K4510 side of the
Tube, any more than a Kaypro could.

It is not shipped, because it is Borland's. Put \texttt{TURBO.COM} and
\texttt{TURBO.MSG} into \pth{fs/CPM/H/3} --- the folder that CP/M's
drive \verb|P:| is a shortcut to --- and install it for a ``VT100'' or
``ANSI'' terminal, which is what JIM is. Then, from the K/OS prompt:

\begin{type}
TURBO
\end{type}

\noindent \verb|TURBO| is an alias that \pth{/SYSTEM/ETC/STARTUP.SAMPLE}
defines, for \verb|CPM K-TURBO|: boot CP/M, go to \verb|P:|, start Turbo,
and --- because the submit ends in \verb|EXIT| --- come all the way back
to the K/OS prompt when you leave it.\shipif{chapter:09-cpm}{ Chapter~\ref{cha:cpm}
explains that arrangement.}

\screen{shots/turbo.png}{Turbo Pascal~3 starting on the Z80, its
terminal drawn by JIM.}

Inside, it is 1985. The editor is WordStar's: Ctrl-E, S, D and X move
(the arrow keys are sent as those, so they work too; Home, End, PgUp and
PgDn are not, and do not), and Ctrl-K then D leaves the editor for the
menu. \verb|R| runs the program in memory;
\verb|O| then \verb|C| makes it compile to a \texttt{.COM} you can run
from the CP/M prompt by name.

\section{Mad Pascal, the cross-compiler}
\emph{Mad Pascal} is the other kind entirely --- a cross-compiler, like
the cc65 this machine's ROM is built with. You write on the desktop,
\texttt{mp} compiles to 6502 assembly, MADS assembles it, and the result
is a \texttt{.prg} beside the source in \pth{/LANG/PASCAL} that the
shell's \verb|RUN| loads like any other program --- and, since 2026-09-08,
one the machine can make for itself: \verb|PAS HELLO| compiles
\texttt{HELLO.PAS} in the directory you are standing in
(Chapter~\ref{cha:linux}). It is Tomasz Biela's compiler (MIT), a
Turbo-Pascal-flavoured language with 8-, 16- and 32-bit integers, fixed
and floating point, strings, records, pointers, units and inline
assembly, made for the Atari and since taught the C64, the X16 and the
Neo6502. The K4510 is its newest target.

\begin{type}
PMANDEL
\end{type}

\screen{shots/pmandel.png}{\texttt{PMANDEL}: the Mandelbrot set in text,
sixteen colours through JIM, in 250 frames.}

That is \texttt{demo/pas/pmandel.pas}: the Mandelbrot set, 78 by 28
cells in the machine's sixteen colours, in a little over four seconds
at the desktop's 40.5~MHz --- fixed point, 32-bit integers, forty lines
of Pascal.
\verb|PSIEVE| is the classic sieve, ten passes timed by the frame
counter; \verb|HELLO| prints, shows the colours, and echoes what you
typed after its name. All of them are already compiled and sitting
beside their sources in \pth{/LANG/PASCAL}, so you can run them without
installing anything.

\subsection*{Building}
On the desktop, with FPC and a Mad-Pascal checkout (the installer
patches its target table and rebuilds \texttt{mp}). It looks for the
checkout in \pth{~/Projects/K4510/toolchain/Mad-Pascal} unless
\texttt{MP\_DIR} says otherwise:

\begin{verbatim}
make pascal-install    # once; MP_DIR=... to point elsewhere
make pascal            # fs/LANG/PASCAL/*.PAS -> *.prg beside them
\end{verbatim}

A new program is a file in \pth{fs/LANG/PASCAL}; \texttt{make} finds
it, and so does \verb|PAS| on the machine itself.
The compiled \texttt{.prg} files are committed, so a machine without
the toolchain still runs them, and the test suite (\texttt{test/pastest.sh})
runs \verb|HELLO| and \verb|PSIEVE| through the ROM.

\subsection*{The target}
Everything a Pascal program needs from the machine is in
\texttt{pascal/} of the repository, laid out as it lives inside a
Mad-Pascal checkout: the runtime base (\texttt{base/rtl6502\_k4510.asm}
and \texttt{base/k4510/}), where \texttt{@putchar} writes to JIM, the
terminal (Chapter~\ref{cha:tube}); the SYSTEM and CRT units' machine
halves (\texttt{lib/*\_k4510.inc}); and a \texttt{k4510} unit. The
consequences for the programmer:

\begin{itemize}
\item \verb|Write| and \verb|WriteLn| go through JIM, so the CRT unit is
  the real thing: \verb|GotoXY|, \verb|TextColor| and
  \verb|TextBackground| (the palette's constants: \verb|BLUE|,
  \verb|YELLOW|, \verb|LIGHT_GREEN|\ldots), \verb|ClrScr|, \verb|ClrEol|,
  \verb|InsLine|/\verb|DelLine|, \verb|WhereX|/\verb|WhereY|,
  \verb|CursorOn|/\verb|CursorOff|, \verb|ReadKey| and \verb|KeyPressed|
  on the keyboard device, \verb|Delay| and \verb|Pause| on the frame
  counter, and \verb|Sound| on the machine's sound sequencer --- which
  changed its contract when the OPL2 became the machine's chip:
  \verb|Sound| takes a channel and a pitch in quarter-semitones and
  holds the note until \verb|NoSound|. \verb|TextMode(0)| puts the
  ROM's screen back.
\item \verb|ParamCount| and \verb|ParamStr| come from the shell
  line (\texttt{RUN HELLO one two}, or just \texttt{HELLO one two});
  \verb|ParamStr(0)| is the whole tail.
\item \verb|uses k4510| gives every chip as a typed variable at its
  address --- \verb|VICKY[reg]|, \verb|DMA_SRC|,
  \verb|FS_CMD|, \verb|SYS_FRAMES|, \verb|MATH_F[]|, \verb|TERM_CX|
  and the rest --- plus \verb|FarPeek|/\verb|FarPoke| into the 256~MB
  through the 45GS02's flat addressing, \verb|DmaCopy|/\verb|DmaFill|,
  \verb|Shell('DIR')|, \verb|LoadFile|/\verb|SaveFile| through the ROM,
  \verb|WaitVBlank|, and \verb|TermWrite| for raw escape sequences ---
  each described under \emph{The k4510 unit}, below.
\item Memory: code from \texttt{\$0800}, the program's own; zero page
  \texttt{\$22--\$3F} for the compiler's registers and
  \texttt{\$64--\$A3} for its expression stack (the ROM keeps
  \texttt{\$02--\$21} and \texttt{\$F0--\$F9}); \texttt{\$0300} the string
  buffer. Programs return to the shell with the ROM's state intact.
\end{itemize}

\subsection*{Floating point on the MATH unit}
\texttt{single} (IEEE-754, 32-bit) does not run in software here. The
runtime's add, subtract, multiply, divide, compare, \verb|Trunc|,
\verb|Round| and \verb|Frac| are the MATH unit at \texttt{\$D700}
(Chapter~\ref{cha:io}): each is a few register moves and one write to
FOP, and the unit answers in the next cycle. \texttt{PFLOAT} times five
thousand rounds of multiply, divide, add and subtract: 5 frames on the
unit, 24 with the software library (assemble with
\verb|mads -d:SOFTFLOAT=1| to get it back for comparison). The
transcendentals are in the \texttt{k4510} unit --- \verb|MathSqrt|,
\verb|MathSin|, \verb|MathCos|, \verb|MathTan|, \verb|MathAtan|,
\verb|MathAtan2|, \verb|MathExp|, \verb|MathLn|, \verb|MathPow|,
\verb|MathFloor| --- one register write each. SYSTEM's own
\verb|Sqrt|/\verb|Sin|/\verb|Cos|/\verb|ArcTan|/\verb|Exp|/\verb|Ln| on
\texttt{single} run on the unit too: the installer patches
\texttt{system.pas} under \verb|{$ifdef k4510}|.

\subsection*{Graphics on VICKY}
\verb|uses graph| draws on VICKY's layer~1: a 640$\times$480 8-bit bitmap
at \texttt{\$200000} --- the one BBC BASIC's ULA uses too --- over the
console, which stays visible wherever the picture is transparent
(colour~0), as in a BBC MODE. \verb|InitGraph(0)| switches the video to
480 lines and clears the bitmap; \verb|CloseGraph| puts the console's
mode back. \verb|PutPixel| and \verb|GetPixel| go through the 45GS02's
flat addressing (the MEGA65 multiplier does the row arithmetic);
\verb|Line|, \verb|LineTo| and the K4510 additions \verb|DrawLine|,
\verb|FillBar| and \verb|FillTriangle| are the blitter's LINE, FILL and
TRIANGLE operations --- one register write each; \verb|Circle|,
\verb|Ellipse|, \verb|Rectangle| and the rest of Mad Pascal's generic
GRAPH build on those; \verb|OutTextXY| writes with the console's font.
\verb|PGRAPH| is the demonstration.

\screen{shots/pgraph.png}{\texttt{PGRAPH}: bars, a star and triangles by
the blitter, circles by pixel, text from the ROM's font --- and the
console's own line still on top.}

\subsection*{The k4510 unit}
\verb|uses k4510| is the machine itself, for a program that wants more
than CRT and GRAPH give it. Most of the unit is the chips' registers as
variables at their addresses (below); the rest is ten routines for memory,
files and the shell, and ten for the MATH unit. \texttt{KUNIT} in
\pth{/LANG/PASCAL} runs each of them once, and is the source of the
examples here.

\paragraph{Far addresses.} The unit's addresses are \texttt{cardinal}s:
physical addresses in the 256\,MB, \texttt{\$0000000} to
\texttt{\$FFFFFFF}. The bottom 64\,KB of it is the CPU's own view, so
\verb|cardinal(@buf)| is the far address of a Pascal variable, and the
DMA routines move between the two. While a program runs, the far memory
from \texttt{\$400000} to \texttt{\$FCFFFFF} --- some 249\,MB --- is
its own. Below that the machine keeps the fonts (\texttt{\$10000}) and
the console's screen (\texttt{\$30000}), and GRAPH draws at
\texttt{\$200000}; above it are the ROM, the sideways ROMs and the
shell's own buffers.

\begin{cmdblock}
\cmdentry{FarPeek(a: cardinal): byte}{far memory}{The byte at far
  address \texttt{a}, by the 45GS02's flat addressing.}
\cmdentry{FarPoke(a: cardinal; v: byte)}{far memory}{Writes \texttt{v}
  there.}
\cmdentry{DmaCopy(src, dst, len: cardinal)}{far memory}{Copies
  \texttt{len} bytes; overlapping is safe. The DMA engine has finished
  before the call returns.}
\cmdentry{DmaFill(dst, len: cardinal; v: byte)}{far memory}{Fills
  \texttt{len} bytes with \texttt{v}.}
\cmdentry{LoadFile(const name: string; dest: cardinal): cardinal}{files}{Reads
  the whole file to far address \texttt{dest} and returns its length;
  0 if it could not (no such file --- or an empty one).}
\cmdentry{SaveFile(const name: string; src, len: cardinal): boolean}{files}{Writes
  \texttt{len} bytes from \texttt{src} to the file, replacing it; true if
  it worked.}
\cmdentry{Shell(const cmd: string)}{the shell}{Runs one line as if typed
  at the prompt: \texttt{DIR}, \texttt{CD}, \texttt{COPY}, \texttt{MKDIR},
  \texttt{ECHO} and the rest. Its output goes to the screen; the result is
  the byte at \texttt{\$03FF} (Chapter~\ref{cha:memory}), 0 if the
  command worked --- read it straight after: it is the last byte of the
  runtime's string buffer, and the next string operation may overwrite it.}
\cmdentry{WaitVBlank}{timing}{Waits for the next frame, a sixtieth of a
  second. \texttt{SYS\_FRAMES} is the frame count's low byte, for timing.}
\cmdentry{TermWrite(const s: string)}{the terminal}{Sends \texttt{s} to
  JIM unchanged, VT100 escape sequences and all.}
\cmdentry{MathSqrt, MathSin, MathCos, MathTan, MathAtan, MathExp, MathLn,
  MathFloor (x: single): single}{MATH unit}{One write to the unit each;
  angles in radians, \texttt{MathLn} the natural log.}
\cmdentry{MathAtan2(y, x: single): single}{MATH unit}{The angle of the
  point (x,\,y), $-\pi$ to $\pi$.}
\cmdentry{MathPow(x, y: single): single}{MATH unit}{$x^y$.}
\end{cmdblock}

\noindent Names and command lines are cut at 95 characters, the length
of the shell's own line; a name is the shell's too, relative to the
current directory or from \pth{/}.

\begin{verbatim}
uses k4510;
const FAR0 = $400000;
var buf: array[0..255] of byte;
    RC: byte absolute $03FF;            { the shell's result }
    len: cardinal;
begin
  DmaFill(FAR0, 256, 0);
  FarPoke(FAR0 + 7, 99);
  DmaCopy(FAR0, cardinal(@buf), 256);   { far -> buf: buf[7] = 99 }
  if SaveFile('/HOME/DATA.BIN', FAR0, 256) then
    len := LoadFile('/HOME/DATA.BIN', FAR0 + $1000);   { 256 }
  Shell('CD /NO-SUCH-DIR');             { RC <> 0: it failed }
end.
\end{verbatim}

\paragraph{The registers.} Every chip is a variable at its address, typed
to fit, under the names of Appendix~\ref{cha:registers}: \verb|VICKY[]|
(the 256 bytes at \texttt{\$D000}), the keyboard (\verb|KBD| pops the next
key, \verb|KBDST| says whether one is waiting and which modifiers are
down), \verb|DMA_SRC|/\verb|DST|/\verb|LEN|/\verb|CMD|, the file
system's \verb|FS_CMD| and its friends (with the \verb|FS_OPEN_READ|\ldots
command constants), \verb|SYS[]| and \verb|SYS_FRAMES|, the MATH unit's
\verb|MATH_F[0..7]| and \verb|MULTINA|/\verb|MULTINB|/\verb|MULTOUT|,
\verb|TUBE[]|, \verb|NET[]|, and JIM's \verb|TERM|, \verb|TERM_CX|,
\verb|TERM_CY|, \verb|TERM_FG|, \verb|TERM_BG| and the rest. Reading or
writing one is a single load or store.

\paragraph{A bank window.} \verb|FarPeek| is a call per byte. For a
block of far memory used as an ordinary array --- indexed, at full speed
--- point an 8\,KB block of the CPU's view at it with a bank register
(Chapter~\ref{cha:memory}) and declare the array \verb|absolute| there.
The unit has no names for the bank registers; declare the one you use:

\begin{verbatim}
var win: array[0..8191] of byte absolute $8000;   { block 4 }
    BANK4: cardinal absolute $D610;               { $D600 + 4*4 }
    BANK4HI: byte absolute $D613;
begin
  BANK4 := $400000;      { open: win[] is $400000-$401FFF }
  win[100] := 1;         { lands in far memory }
  BANK4HI := $FF;        { shut: $8000 is the program's RAM again }
end.
\end{verbatim}

\noindent Writing the whole \texttt{cardinal} switches the block on as
well (its top byte has bit~7 clear); \texttt{\$FF} in that byte switches
it off. Block $n$ is \texttt{\$2000}$\times n$ to
\texttt{\$2000}$\times n +$\texttt{\$1FFF}. Use one the program does
not occupy --- its code and variables run up from \texttt{\$0800}
(\texttt{BANKWIN}'s end at \texttt{\$137A}, so blocks 1 to 4 are
free) --- and never 0 (zero page and stack) or 5 to 7, which the ROM
banks in and out around every system call. \texttt{BANKWIN} is the
full demonstration, timed against \verb|FarPeek|: the window sums
64\,KB in 8 frames, \verb|FarPeek| in 17.

\paragraph{The rules.}
\begin{itemize}
\item \textbf{Code stays in the 64\,KB}; data does not have to. A
  program is at most \texttt{\$0800}--\texttt{\$CFFF}, 50\,KB with its
  variables; everything else goes in far memory.
\item \textbf{Do not \texttt{Shell} a program.} \verb|RUN| loads the
  other program over yours (Pascal programs load at \texttt{\$0800},
  C programs at \texttt{\$6000}) and there is nothing to come back to.
  The shell's own commands are fine.
\item \textbf{Inline assembly that borrows \texttt{:bp}} must put its
  low byte back to 0. Mad Pascal indexes any array over 256 bytes
  through \verb|(:bp),y| and counts on it; the unit's own routines
  restore it.
\end{itemize}
